We consider a sequential clinical trial with $K$ treatments, also called arms. Patients arrive sequentially and are assigned to treatments adaptively. In the following we denote \([K]:= \left\{1, \dots, K\right\}\).

\paragraph*{Arm's distributions}
For each patient $m$, let \(X_m = (X_{m,1}, \ldots, X_{m,K})\) denote the assignment vector, where $X_{m,k}=1$ indicates that the patient \(m\) is assigned to arm \(k \in [K]\) and $X_{m,k}=0$ otherwise, thus we have $\sum_{k=1}^{K} X_{m,k} = 1$.

Consider the sequence of \textit{i.i.d.} random variables \(\left\{\xi_{m,k}, m=1, \dots\right\}\)  where \(\xi_{m,k}\) is the response of the \(m^{th}\) patient on treatment \(k\). For a treatment $k \in [K]$, the parameter of interest in this paper is the expectation of these responses, we denote it as \(\theta_k := \mathbb{E}\left[\xi_{1,k}\right]\). We also define the vector of these parameters, \(\Theta = (\theta_1, \ldots, \theta_K) \in \mathbb{R}^K\).

The total number of patients assigned to each treatment up to time $n$ is denoted by \(N_n = (N_{n,1}, \ldots, N_{n,K})\) where \(N_{n,k} = \sum_{j=1}^n X_{j,k}\).
The normalized vector $N_n/n$ represents the observed allocation proportions at time $n$. Controlling the long-term behavior of these proportions is the central objective of adaptive targeting mechanisms.

\begin{remark}
Prior work often consider a more general setting in which both the observation $\xi_{m,k}$ %can be a $d$-dimensional vector , i.e. $\xi_{m,k} = (\xi_{m,k,1}, \ldots, \xi_{m,k,d}) \in \mathbb{R}^d$,
and the parameter of the distribution $\theta_k$ can be multi-dimensional (e.g., $\theta_k$ is the mean and covariance of a Gaussian arm). %For example, when treatment are modelled by Gaussian distributions with unknown mean and variance, $\xi_{m,k}$ is one-dimensional, while $\theta_k$ is two-dimensional.
For the sake of clarity, we focus our investigation on the (doubly) uni-dimensional case, for distribution parameterized by their means. This encompasses the example of arms that belong to a one-dimensional exponential family, that has received particular attention in the multi-armed bandit literature \citep{KLUCBJournal,garivier2016optimal}. However,  our analysis and results extend to these more complex settings without substantial modification.
\end{remark}

\paragraph*{Response-Adaptive Randomization (RAR) procedures}

Let $\mathcal{F}_m = \sigma\!\left((X_i,\xi_i)_{1 \le i \le m}\right)$ 
denote the filtration generated by the treatment assignments and observed 
outcomes of the first $m$ patients. A \textit{RAR} procedure is a sequence
$(p_m)_{m\ge1}$ of $\mathcal{F}_m$-measurable random vectors taking values 
in the probability simplex
\[
\Delta_K := \left\{ p \in [0,1]^K : \sum_{k=1}^K p_k = 1 \right\}.
\]
At step $m+1$, the treatment assignment $X_{m+1}$ is drawn conditionally 
on $\mathcal{F}_m$ according to $p_m$, that is,
\[
\mathbb{P}(X_{m+1} = k \mid \mathcal{F}_m) = p_{m,k},
\qquad k = 1,\dots,K.
\]


\paragraph*{Adaptive Targeting Mechanisms} Given some target function $\rho : \mathcal{H} \rightarrow \Delta_K$, i.e. a function such that, for any $z \in \mathcal{H}$,
\[
\rho(z) = (\rho_1(z), \ldots, \rho_K(z))
\]
satisfies $\sum_{k=1}^{K}\rho_{k}(z)=1$, we propose in this work RAR procedures under which the empirical allocation $N_n/n$ converges to the target allocation $v := \rho(\Theta)$. We refer to these procedures as adaptive targeting mechanisms.

Different choices of $\rho$ correspond to different design objectives, such as equal allocation, Neyman allocation for power, or ethically motivated rules \citep{pin2024response}.
Existing adaptive targeting mechanisms include DBCD \citep{hu2004asymptotic} and ERADE \citep{hu2006theory}, which we extend in the next section to $K \geq 2$ arms.


\paragraph*{Notation}
In the paper we will use the notations $o_p(\cdot)$ and $O_p(\cdot)$, which are analogues of the deterministic little-$o$ and big-$O$ notations for convergence in probability.
More formally, for a sequence $\{X_n\}$ of random variables and $\left(a_n\right)_{n \geq 1}$ a sequence of positive real numbers.
We say that $X_n = o_p(a_n)$ if for all $\varepsilon > 0$, \(\lim_{n \to \infty} \mathbb{P}\!\left( \frac{|X_n|}{a_n} > \varepsilon \right) = 0.\)
We say that $X_n = O_p(a_n)$ if for all $\varepsilon > 0$, there exists $M > 0$ such that \(\limsup_{n \to \infty} \mathbb{P}\!\left( \frac{|X_n|}{a_n} > M \right) \leq \varepsilon.\)		
	
